\documentclass[10pt,a4paper]{amsart}
\usepackage[margin=19mm]{geometry}
\usepackage{amsmath,amssymb,mathtools}
\usepackage[scr=rsfs,cal=euler]{mathalfa}
\usepackage{xfrac}
\usepackage[unicode,hidelinks,bookmarks=false]{hyperref}
\newcommand{\ie}{i.e.\ }
\newcommand{\cf}{cf.\ }
\newcommand{\resp}{respectively\ }
\newcommand{\Subsection}{\subsection}
\providecommand{\nobreakdash}{\nobreak\hbox{-}}
\setlength{\emergencystretch}{2em}
\multlinegap0pt
\usepackage{amssymb}
\usepackage{csquotes}
\newtheorem{theorem}{Theorem}
\title{Rough Weighted Ideal Convergence and Korovkin-Type Approximation via weighted equi-ideal convergence}
\author{Tamim Aziz, Sanjoy Ghosal}
\date{Source: arXiv:2512.18676v1 (CC BY 4.0)}
\begin{document}
\maketitle

\noindent\emph{Source and licence.} Rough Weighted Ideal Convergence and Korovkin-Type Approximation via weighted equi-ideal convergence. Version arXiv:2512.18676v1 (CC BY 4.0), distributed by arXiv under the Creative Commons Attribution 4.0 International licence, http://creativecommons.org/licenses/by/4.0/, as recorded in the arXiv licence metadata for this exact version and shown on the abstract page https://arxiv.org/abs/2512.18676v1. Full licence notice: \texttt{licenses/arxiv-2512.18676-CC-BY-4.0.txt}.\par\smallskip

\noindent\textit{Editorial scope.} Counted unit: the article's theorem korokvingequiideal (source0.tex lines 1189--1199). The article's complete original proof of it is delivered with it (source0.tex lines 1200--1275, one \texttt{proof} environment of 76 lines). No mathematics is written, repaired or re-derived: the statement and the proof are the article's own lines, in the article's own notation, and no retained line is substituted or re-flowed.  No further source-local context is needed: every cross-reference of every retained line resolves inside the two delivered spans.  Nothing was omitted from any retained span, so the package carries no omission record.  No retained line contains a citation, so the wrapper carries no bibliography and no bibliography entry.  No retained line contains a figure environment, an image-inclusion command or drawing code, so no image is delivered and the package declares no asset dependency.  Editorial formatting only, in addition: an AMS \texttt{amsart} preamble that restates the article's own declarations and macros used by the retained lines, namely the environments \texttt{theorem} on separate counters, together with the standard packages those lines need.  The retained environments print in this document as Theorem 1 in order of appearance, whereas the article prints the same items with the numbering of its own full document.  The complete native source of the article as distributed in the arXiv e-print for this exact version is bundled unchanged as \texttt{sources/191411/source0.tex}.
	\begin{theorem}	\label{korokvingequiideal}
		Let $\mathcal{I}_{\varphi}$ be an analytic $P$-ideal, and let $\{\omega_t\}_{t \in \mathbb{N}}$ be a weighted sequence that is $\mathcal{I}_{\varphi}$-bounded. Consider a compact subset $K \subset \mathbb{R}$, and let $\{L_t : C(K) \to C(K)\}_{t \in \mathbb{N}}$ be a sequence of positive linear operators.
		 Then, for all $f \in C(K)$, 
	\begin{equation*}
	L_t(f) \to f \quad \text{($\omega-equi-\mathcal{I}_{\varphi}$)} \quad \text{on } K \tag{a}
\end{equation*}
	if and only if
\begin{equation}
		L_t(e_i) \to e_i \quad \text{($\omega-equi-\mathcal{I}_{\varphi}$)} \quad \text{on } K \quad \text{with } e_i(x) = x^i, \ i = 0, 1, 2\tag{b}.
\end{equation}
\end{theorem}

	\begin{proof}
		The forward implication (a) $\Rightarrow$ (b) is immediate, since each monomial $e_i \in C(X)$ for $i = 0, 1, 2$.
	

To establish the converse, assume that condition (b) holds. Since the weighted sequence $\{\omega_{t}\}_{t\in\mathbb{N}}$ is $\mathcal{I}_{\varphi}$-bounded, there exists $\mu>0$ such that 
$$
A=\{t\in\mathbb{N}:\omega_t>\mu\}\in\mathcal{I}_{\varphi}.
$$
Now, fix an arbitrary function $f \in C(K)$ and a point $c \in K$. By the continuity of $f$ at $c$, for any $\varepsilon > 0$, there exists a $\delta > 0$ such that
\[
|f(x) - f(c)| \leq \frac{\varepsilon}{\mu} \quad \text{whenever } |x - c| < \delta.
\]

Define the set $K_\delta := [c - \delta, c + \delta] \cap K$ and denote by $\chi_D$ the characteristic function of a set $D$. Then we can decompose
\[
|f(x) - f(c)| = |f(x) - f(c)| \chi_{K_\delta}(x) + |f(x) - f(c)| \chi_{K \setminus K_\delta}(x).
\]

This leads to the bound
\[
|f(x) - f(c)| \leq \frac{\varepsilon}{\mu} + 2B \frac{(x - c)^2}{\delta^2} \quad \text{for all } x \in K,
\]
where $B := \|f\|_{C(K)}$.

Using the positivity and linearity of the operators $L_t$, we obtain
\[
|L_t(f; c) - f(c)| \leq \frac{\varepsilon}{\mu} + (\frac{\varepsilon}{\mu} + B) |L_t(e_0; c) - e_0(c)| + \frac{2B}{\delta^2} L_t\left((x - c)^2; c\right).
\]

Expanding the quadratic term,
\[
(x - c)^2 = x^2 - 2cx + c^2 = e_2(x) - 2c e_1(x) + c^2 e_0(x),
\]
and applying $L_t$ yields
\[
L_t\left((x - c)^2; c\right) = L_t(e_2; c) - 2c L_t(e_1; c) + c^2 L_t(e_0; c).
\]

Consequently,
\begin{align*}
|L_t(f; c) - f(c)| \leq 
&\frac{\varepsilon}{\mu}+ \left(\frac{\varepsilon}{\mu} + B + \frac{2B c^2}{\delta^2} \right) |L_t(e_0; c) - e_0(c)|\\& 
+ \frac{4B c}{\delta^2} |L_t(e_1; c) - e_1(c)| + \frac{2B}{\delta^2} |L_t(e_2; c) - e_2(c)|.
\end{align*}

Thus, we can write
\begin{equation*}
	|L_t(f; c) - f(c)| \leq \frac{\varepsilon}{\mu} + B' \sum_{i=0}^2 |L_t(e_i; c) - e_i(c)|, 
\end{equation*}
where
\[
B' := \frac{\varepsilon}{\mu} + B + \frac{2B}{\delta^2} \left( \|e_2\|_{C(K)} + 2\|e_1\|_{C(K)} + 1 \right).
\]
This ensures that
$$
\omega_t |L_t(f; c) - f(c)| \leq \varepsilon + B' \sum_{i=0}^2 \omega_t|L_t(e_i; c) - e_i(c)|~\mbox{ for each}~n\in\mathbb{N}\setminus A. 
$$
Let us set 
\begin{align*}
&\Psi_j(c,\varepsilon)=\varphi(\{t\in\mathbb{N}\setminus A:\omega_t |L_t(f; c) - f(c)|>\frac{\varepsilon}{2}\}\setminus [1,j])\\
&\Psi_{i,j}(c,\varepsilon)=\varphi(\{t\in\mathbb{N}\setminus A:\omega_t |L_t(e_i; c) - e_i(c)|>\frac{\varepsilon}{6B'}\}\setminus [1,j])~\mbox{ for }~i=0,1,2.
\end{align*}
Since $c\in K$ was chosen arbitrarily and  $\varphi$  is monotone non-decreasing, we obtain that
\begin{align*}
\Psi_j(x,\varepsilon)\leq \displaystyle \sum_{i=0}^{2}\Psi_{i,j}(x,\varepsilon)~\mbox{for all $x\in K$ and for each }~j\in\mathbb{N}.
\end{align*}
By the hypothesis, it follows that the sequence of functions $\{\Psi_{i,j}(x,\varepsilon)\}_{j\in\mathbb{N}}$ is uniformly convergent to zero function on $K$ for each $i=0,1,2$. As a consequence, $\{\Psi_j(x,\varepsilon)\}_{j\in\mathbb{N}}$ converges uniformly to zero function on $K$. It now remains to show that the sequence of functions $\{h_j(x,\varepsilon)\}_{j\in\mathbb{N}}$, defined by 
$$
h_j(x,\varepsilon)=\varphi(\{t\in\mathbb{N}:\omega_t |L_t(f; x) - f(x)|>\frac{\varepsilon}{2}\}\setminus [1,j]),
$$
also converges uniformly to zero function on $K$. This holds since
$$
0\leq h_j(x,\varepsilon)\leq \Psi_j(x,\varepsilon)+\varphi(A\setminus [1,j])~\mbox{ and }~A\in\mathcal{I}_{\varphi}.
$$
Hence the implication (b) $\Rightarrow$ (a) holds as well.
\end{proof}

\end{document}
